\subsection{2D Inviscid Bump}
\label{SEC:Testcases:Bump}

\begin{figure}[h!tb]
  \begin{center}
    \insertimage[width=6.0in]{test_cases/bump}
    \caption[Inviscid Gaussian bump test case.]%
            {Inviscid Gaussian bump test case, \PO[3] solution.}
    \label{FIG:Bump}
  \end{center}
\end{figure}

This test case involves the subsonic flow over a smooth Gaussian bump in a
channel, and is aimed at testing high-order methods for the computation of
internal flow with a high-order curved boundary representation.
%
In this subsonic flow problem, the geometry is smooth and symmetric in $x$,
and so should the flow properties.

This test case can be used to experiment with how the grid order affects
solution accuracy for various solution orders.

All of the provided grid files contain the same number of grid cells (256),
and the names of the grid files refer to the order of the cells contained
in the grid.
%
The 'bump.P1.msh' grid contains linear    (\PO[1]) cells,
the 'bump.P2.msh' grid contains quadratic (\PO[2]) cells,
the 'bump.P3.msh' grid contains cubic     (\PO[3]) cells, and
the 'bump.P4.msh' grid contains quartic   (\PO[4]) cells.
%
%This case reaches convergence in just over 10 minutes running in serial,
%and roughly 60 seconds using 16 MPI processes.
%
\Fref{FIG:Bump} shows the result on the \PO[3] grid.

\SeeAlso{
\inp[nmlInitialization::Initialization]{Initialization}=3
}

